\chapter{The Four Horsemen: Penjaga dari Neraka}

\begin{insightbox}
\textit{``Ada empat perilaku yang bisa memprediksi perceraian dengan akurasi 91\%. Mereka adalah Four Horsemen of the Apocalypse untuk pernikahan: Criticism, Contempt, Defensiveness, dan Stonewalling.''}
\end{insightbox}

\section{Pengantar: Prediktor Perceraian}

Dr. John Gottman menemukan bahwa \textbf{ Four Horsemen} (Empat Jahanam) adalah penanda terkuat kehancuran hubungan. Keempatnya adalah:

\begin{center}
\begin{tikzpicture}[
    horse/.style={draw=#1, line width=3pt, rounded corners=15pt,
                  fill=#1!15, minimum width=12cm, minimum height=2.5cm,
                  align=left, text width=11cm}
]
    \node[horse=accent] (critic) at (0,6) {
        \textbf{\large 1. CRITICISM (Kritik)}\\[0.2cm]
        Menyerang \textbf{karakter/pribadi} pasangan, bukan perilaku spesifik
    };
    
    \node[horse=secondary] (contempt) at (0,3) {
        \textbf{\large 2. CONTEMPT (Penghinaan)}\\[0.2cm]
        Perasaan \textbf{superior} disertai sarkasme, ejekan, mata melotot
    };
    
    \node[horse=warmgray] (defens) at (0,0) {
        \textbf{\large 3. DEFENSIVENESS (Bela Diri)}\\[0.2cm]
        Membela diri, membalikkan tuduhan, memainkan peran korban
    };
    
    \node[horse=primary] (stone) at (0,-3) {
        \textbf{\large 4. STONEWALLING (Membisu)}\\[0.2cm]
        Menarik diri, tidak merespons, membangun ``tembok batu''
    };
    
    % Arrows showing progression
    \draw[->, line width=2pt, accent!50] (critic.south) -- (contempt.north);
    \draw[->, line width=2pt, secondary!50] (contempt.south) -- (defens.north);
    \draw[->, line width=2pt, warmgray!50] (defens.south) -- (stone.north);
    
    % Label
    \node[right=0.5cm of critic, font=\small, accent] {Awal};
    \node[right=0.5cm of stone, font=\small, primary] {Akhir};
\end{tikzpicture}
\end{center}

\section{Horseman 1: Criticism (Kritik)}

\subsection{Bedakan: Complaint vs Criticism}

\begin{center}
\renewcommand{\arraystretch}{1.5}
\begin{tabular}{|>{\raggedright}p{6cm}|>{\raggedright\arraybackslash}p{9cm}|}
\hline
\rowcolor{sagegreen!20}
\textbf{COMPLAINT (Benar)} & \textbf{CRITICISM (Salah)} \\
\hline
\textit{``Aku marah kamu nggak sapu dapur. Kita sepakat gantian. Bisa tolong sekarang?''} & \textit{``Kenapa kamu selalu lupa? Aku benci harus selalu sapu. Kamu nggak peduli.''} \\
\hline
Fokus pada \textbf{perilaku spesifik} & Menyerang \textbf{karakter} (``selalu'', ``kamu nggak peduli'') \\
\hline
Ada 3 komponen: Perasaan + Situasi + Request & Global, menyeluruh, menghakimi \\
\hline
\end{tabular}
\end{center}

\subsection{Formula Complaint yang Benar}

\begin{praktikbox}[Struktur Complaint Sehat]
\textbf{``Aku [perasaan] ketika [situasi spesifik]. Aku butuh [request].''}

\textbf{Contoh:}
\begin{itemize}
    \item ``Aku \textbf{marah} ketika \textbf{kamu nggak sapu dapur semalam}. Aku butuh \textbf{kamu sapu sekarang}.''
    \item ``Aku \textbf{kecewa} ketika \textbf{kamu nggabarin aku kalo telat}. Aku butuh \textbf{kabarin minimal 30 menit sebelumnya}.''
    \item ``Aku \textbf{merasa diabaikan} ketika \textbf{kamu main HP saat aku bicara}. Aku butuh \textbf{kamu taruh HP dan dengarkan dulu 10 menit}.''
\end{itemize}
\end{praktikbox}

\subsection{Kata-kata Pembunuh}

Hindari kata-kata ini dalam konflik:
\begin{itemize}
    \item ``Kamu \textbf{selalu}...''
    \item ``Kamu \textbf{tidak pernah}...''
    \item ``Kamu \textbf{terlalu}...''
    \item ``\textbf{Kenapa} kamu...'' (terdengar seperti kritik)
    \item ``Apa yang \textbf{salah} denganmu?''
\end{itemize}

\section{Horseman 2: Contempt (Penghinaan)}

\subsection{Ini adalah yang Paling Berbahaya}

\begin{warningbox}[Contempt = Poison]
Contempt adalah \textbf{prediktor tunggal terkuat} perceraian. Ketika ada contempt, hubungan dalam bahaya serius.

\textbf{Tanda-tanda Contempt:}
\begin{itemize}
    \item Sarkasme, ejekan, name-calling
    \item Mata melotot, sneering (menyeringai menghina)
    \item Body language meremehkan (menyilang tangan, mata ke langit)
    \item Menyindir cara pasangan bicara/berpakaian/berperilaku
    \item Merasa ``aku lebih baik darimu''
\end{itemize}
\end{warningbox}

\subsection{Contoh Perilaku Contempt}

\begin{center}
\renewcommand{\arraystretch}{1.4}
\begin{tabular}{|>{\raggedright}p{6cm}|>{\raggedright\arraybackslash}p{9cm}|}
\hline
\rowcolor{accent!20}
\textbf{Konteks} & \textbf{Contempt dalam Aksi} \\
\hline
Uang & ``Kamu itu \textbf{manja} dan \textbf{tidak tahu bersyukur}. Aku berasal dari keluarga yang \textbf{tahu nilai uang}, beda sama kamu.'' \\
\hline
Kerja rumah & ``Ah, \textbf{gitu aja nggak bisa}. Dasar \textbf{bodoh} kamu ini.'' \\
\hline
Pekerjaan & ``\textbf{Ngabisin waktu} aja kerjaanmu itu. \textbf{Siapa sih yang peduli}.'' \\
\hline
\end{tabular}
\end{center}

\subsection{Antidot: Fondness and Admiration}

Satu-satunya antidot untuk contempt adalah \keyterm{Fondness and Admiration System}:

\begin{tipbox}[Cara Membangun Fondness]
\begin{enumerate}
    \item \textbf{Scan for positives:} Sengaja cari hal positif tentang pasangan
    \item \textbf{Express gratitude:} Ucapkan terima kasih setiap hari
    \item \textbf{Reminisce:} Ingat kenangan indah masa lalu
    \item \textbf{Pride:} Ungkapkan kebanggaan pada pasangan
    \item \textbf{Affection:} Tunjukkan kasih sayang secara teratur
\end{enumerate}
\end{tipbox}

\section{Horseman 3: Defensiveness (Bela Diri)}

\subsection{Bentuk-bentuk Defensiveness}

\begin{enumerate}
    \item \textbf{Innocent Victim:} ``Kenapa kamu selalu nyalahin aku?''
    \item \textbf{Reverse Blame:} ``Iya, tapi kamu juga...''
    \item \textbf{Excuse:} ``Aku nggak bisa karena...''
    \item \textbf{Whining:} ``Apa pun yang kulakukan nggak pernah cukup!''
\end{enumerate}

\subsection{Contoh Defensiveness}

\begin{praktikbox}[Dialog Defensiveness]
\textbf{Pasangan:} ``Kamu telat lagi hari ini.''

\textbf{Defensiveness:}
\begin{itemize}
    \item ``Ya iyalah, macet! Bukan salahku!'' (excuse)
    \item ``Kamu aja kemarin telat, kok nggak bilang?'' (reverse blame)
    \item ``Kenapa sih kamu selalu komplain mulu?!'' (innocent victim)
\end{itemize}

\textbf{Non-Defensiveness (Benar):}
\begin{itemize}
    \item ``Iya, maaf. Ada kecelakaan di tol. Aku harusnya kabarin.''
    \item ``Kamu kecewa ya? Maaf bikin kamu khawatir.''
\end{itemize}
\end{praktikbox}

\section{Horseman 4: Stonewalling (Membisu)}

\subsection{Apa itu Stonewalling?}

\keyterm{Stonewalling} adalah ketika seseorang menarik diri dari interaksi, tidak merespons, dan membangun ``tembok batu''.

\begin{warningbox}[Tanda-tanda Stonewalling]
\begin{itemize}
    \item Tidak ada kontak mata
    \item Tidak ada respons verbal (``uh-huh'', ``hmm'')
    \item Body language tertutup (memalingkan tubuh)
    \item Ekspresi wajah datar (``poker face'')
    \item Fisik meninggalkan ruangan
\end{itemize}
\end{warningbox}

\subsection{Kenapa Orang Stonewall?}

\begin{praktikbox}[Penyebab Stonewalling]
\textbf{85\% stonewaller adalah pria} --- bukan karena pria ``jelek'', tapi karena:
\begin{enumerate}
    \item \textbf{Fisiologis:} Pria lebih mudah ``flooded'' (sistem kardiovaskular lebih reaktif)
    \item \textbf{Recovery:} Pria butuh waktu lebih lama untuk pulih dari konflik
    \item \textbf{Protection:} Stonewalling adalah self-protection dari flooding
    \item \textbf{Learned behavior:} Dipetik dari pola orang tua
\end{enumerate}
\end{praktikbox}

\section{Flooding: Akar dari Stonewalling}

\subsection{Apa itu Flooding?}

\keyterm{Flooding} adalah kondisi fisiologis di mana tubuh dalam mode ``fight or flight'':

\begin{center}
\begin{tikzpicture}[
    symptom/.style={draw=accent, rounded corners=8pt, fill=accent!10,
                    minimum width=4cm, minimum height=1.5cm, align=center}
]
    \node[font=\Large\bfseries, accent] at (0,4) {TANDA-TANDA FLOODING};
    
    \node[symptom] at (-4,2) {Jantung berdebar\\$>$100 bpm};
    \node[symptom] at (0,2) {Keringat dingin};
    \node[symptom] at (4,2) {Napas cepat/dangkal};
    \node[symptom] at (-4,0) {Otot tegang};
    \node[symptom] at (0,0) {Sulit berpikir};
    \node[symptom] at (4,0) {Ingin lari/menyerang};
    \node[symptom] at (-2,-2) {Telinga panas};
    \node[symptom] at (2,-2) {Pandangan sempit};
\end{tikzpicture}
\end{center}

\section{SAFE WORD/EMOJI SYSTEM (Prioritas Tinggi)}

\begin{insightbox}
\textbf{Ini adalah tools paling penting dalam bab ini.} Safe word/emoji memungkinkan kalian berkomunikasi saat kata-kata sudah tidak bisa digunakan.
\end{insightbox}

\subsection{Konsep Safe Word}

\begin{praktikbox}[Apa itu Safe Word?]
\textbf{Safe word} adalah kata atau sinyal yang disepakati untuk:
\begin{enumerate}
    \item Menyampaikan ``Aku flooded dan butuh break'' tanpa menjelaskan panjang lebar
    \item Menghentikan spiral negatif \textbf{sebelum} terlalu parah
    \item Memberi ruang untuk self-soothe
    \item Menjaga respect dan safety dalam konflik
\end{enumerate}
\end{praktikbox}

\subsection{Jenis-jenis Safe Signal}

\begin{center}
\renewcommand{\arraystretch}{1.5}
\begin{tabular}{|>{\raggedright}p{3cm}|>{\raggedright}p{5cm}|>{\raggedright\arraybackslash}p{7cm}|}
\hline
\rowcolor{primary!20}
\textbf{Tipe} & \textbf{Contoh} & \textbf{Kapan Digunakan} \\
\hline
\textbf{Safe Word} & ``Pineapple'', ``Kuning'', ``Time-out'' & Saat mulai merasa flooded \\
\hline
\textbf{Safe Emoji} & \emoji{\faPause}, \emoji{\faHandPaper}, \emoji{\faLifeRing} & Saat chatting/texting \\
\hline
\textbf{Safe Gesture} & Mengangkat tangan, membuat ``T'' & Saat sudah terlalu emosional untuk bicara \\
\hline
\textbf{Safe Object} & Benda kecil khusus yang diletakkan & Signal ``Aku butuh waktu'' \\
\hline
\end{tabular}
\end{center}

\subsection{Contoh Safe Word/Emoji System}

\begin{exercisebox}[Membuat Safe System Bersama]
\textbf{Langkah 1: Pilih Safe Word}
\begin{itemize}
    \item Pilih kata yang \textbf{tidak biasa} dalam percakapan kalian
    \item Hindari kata yang mungkin muncul di tengah konflik
    \item Contoh baik: ``Pineapple'', ``Blue moon'', ``Kode merah''
    \item Contoh buruk: ``Stop'', ``Cukup'', ``Sebentar''
\end{itemize}

\textbf{Langkah 2: Pilih Safe Emoji}
\begin{itemize}
    \item Saat chatting, gunakan emoji spesifik
    \item Contoh: \emoji{\faPause} = ``Aku overwhelmed, pause dulu''
    \item Contoh: \emoji{\faLifeRing} = ``Aku flooded butuh break''
    \item Contoh: \emoji{\faClock} = ``Aku butuh 20 menit, nanti balik''
\end{itemize}

\textbf{Langkah 3: Tetapkan Protokol}
\begin{itemize}
    \item Ketika safe word digunakan, \textbf{konflik langsung stop}
    \item Tidak ada ``tambahan kalimat terakhir''
    \item Kedua pihak pergi ke ruang terpisah
    \item Setelah 20-30 menit, kembali dan coba lagi
\end{itemize}
\end{exercisebox}

\subsection{Protokol Safe Word}

\begin{tipbox}[Cara Menggunakan Safe Word dengan Benar]
\textbf{Ketika kamu menggunakan safe word:}
\begin{enumerate}
    \item Katakan dengan jelas (``Aku pakai safe word: Pineapple'')
    \item Jelaskan secara singkat kenapa (opsional tapi bagus): ``Aku flooded, jantungku berdebar''
    \item Katakan kapan kamu akan kembali: ``Aku butuh 20 menit''
    \item Pergi ke ruang lain
    \item Lakukan self-soothing
    \item Kembali tepat waktu yang dijanjikan
\end{enumerate}

\textbf{Ketika pasangan menggunakan safe word:}
\begin{enumerate}
    \item \textbf{Respek segera} --- stop bicara meski kamu ``hampir selesai''
    \item Jangan kejar (don't pursue)
    \item Jangan anggap personal (``Dia nggak peduli pendapatku'')
    \item Manfaatkan waktu untuk self-soothe juga
    \item Siap mendengar saat dia kembali
\end{enumerate}
\end{tipbox}

\subsection{Self-Soothing saat Break}

\begin{praktikbox}[Aktivitas Self-Soothing yang Efektif]
\textbf{Fisiologis:}
\begin{itemize}
    \item Deep breathing (4-7-8)
    \item Progressive muscle relaxation
    \item Mandi air hangat
    \item Jalan-jalan ringan
    \item Minum air hangat
\end{itemize}

\textbf{Mental:}
\begin{itemize}
    \item Meditasi atau mindfulness
    \item Baca buku ringan
    \item Dengarkan musik menenangkan
    \item Tulis jurnal perasaan
\end{itemize}

\textbf{Yang TIDAK boleh dilakukan:}
\begin{itemize}
    \item Ruminate (muter-muter pikiran negatif)
    \item Stalking sosial media pasangan
    \item Hubungi teman untuk ``curhat'' (ini malah memperkuat negativity)
    \item Alkohol atau zat
\end{itemize}
\end{praktikbox}

\section{Antidot Lengkap untuk Four Horsemen}

\begin{center}
\renewcommand{\arraystretch}{1.5}
\begin{tabular}{|>{\raggedright}p{3cm}|>{\raggedright}p{4.5cm}|>{\raggedright\arraybackslash}p{7.5cm}|}
\hline
\rowcolor{primary!20}
\textbf{Horseman} & \textbf{Antidot} & \textbf{Contoh Praktis} \\
\hline
\textbf{Criticism} & \textbf{Gentle Start-up} & ``Aku marah karena...'' (bukan ``Kamu selalu...'') \\
\hline
\textbf{Contempt} & \textbf{Fondness \& Admiration} & Sebutkan 3 hal yang kamu hargai tentang pasangan setiap hari \\
\hline
\textbf{Defensiveness} & \textbf{Take Responsibility} & ``Sebagian ini salahku...'' \\
\hline
\textbf{Stonewalling} & \textbf{Safe Word + Self-Soothe} & ``Aku pakai safe word. Aku butuh 20 menit.'' \\
\hline
\end{tabular}
\end{center}

\section{Action Steps}

\begin{exercisebox}[To-Do: Setup Safe System]
\textbf{Hari ini:} Diskusikan dengan pasangan
\begin{enumerate}
    \item ``Apakah kita pernah mengalami flooding? Kapan?''
    \item ``Bagaimana biasanya kita bereaksi saat flooded?''
    \item ``Apakah ada Four Horsemen yang sering muncul dalam konflik kita?''
\end{enumerate}

\textbf{Buat kesepakatan:}
\begin{enumerate}
    \item Pilih safe word bersama
    \item Pilih safe emoji untuk chat
    \item Tulis protokol dan tempel di kulkas
    \item Latihan: Role-play menggunakan safe word
\end{enumerate}

\textbf{Protokol yang ditulis:}
\begin{itemize}
    \item Safe word kami: \rule{5cm}{0.4pt}
    \item Safe emoji: \rule{3cm}{0.4pt}
    \item Waktu break: \rule{3cm}{0.4pt} menit
    \item Aktivitas self-soothe: \rule{8cm}{0.4pt}
\end{itemize}
\end{exercisebox}
